% Related lecture: SC4001-L08, 2026-10-09
% Source: 8_rnn.pdf, file pp. 18--22 and 56.
% Prepared from templates/lecture-example.tex; no conversational checks.

\exercisegroup{SC4001-L08-EX}{第 08 讲：Recurrent Neural Networks}

\exercisemeta
  {SC4001-L08-EX}
  {2026-10-09}
  {8\_rnn，Example 1（pp.~18--22）、Example 2（p.~56）}
  {前向数值已独立复算；序列实验仅分析题意与依赖结构，未运行训练}
  {课堂学习已完成；2026-10-09 整理}

\exercisequestion{SC4001-L08-E01}{两输入、三隐藏单元 RNN 的四步前向计算}

\textbf{对应知识点：}\relatedtopic{SC4001-L08-T01}{前向计算与参数共享}。

\textbf{题目（pp.~18--22，Example 1）。}隐藏反馈 RNN 有两个输入、三个隐藏单元和两个输出；$h_0=0$，隐藏层使用 $\tanh$，输出层使用逐元素 sigmoid。参数为
\[
 U=\begin{pmatrix}-1&0.5&0.2\\0.5&0.1&-2\end{pmatrix},\qquad
 W=\begin{pmatrix}2&1.3&-1\\1.5&0&-0.5\\-0.2&1.5&0.4\end{pmatrix},
\]
\[
 V=\begin{pmatrix}2&-1\\-1.5&0.5\\0.2&0.8\end{pmatrix},\qquad
 b=\begin{pmatrix}0.2\\0.2\\0.2\end{pmatrix},\qquad
 b_y=\begin{pmatrix}0.5\\0.1\end{pmatrix}.
\]
其中 $b_y$ 即原题的输出 bias $c$。输入依次为
\[
 x_1=(1,2)^\top,\quad x_2=(-1,1)^\top,\quad
 x_3=(0,3)^\top,\quad x_4=(2,-1)^\top.
\]
求各步输出 $y_t$。

\begin{workedsolution}
每步按 $z_t=U^\top x_t+W^\top h_{t-1}+b$、$h_t=\tanh z_t$、$y_t=\sigma(V^\top h_t+b_y)$ 计算，不能把输出 $y_t$ 当成下一步的隐藏状态。

第一步循环项为零：
\[
 z_1=\begin{pmatrix}-1+1+0.2\\0.5+0.2+0.2\\0.2-4+0.2\end{pmatrix}
    =\begin{pmatrix}0.2\\0.9\\-3.6\end{pmatrix},\qquad
 h_1\approx\begin{pmatrix}0.197375\\0.716298\\-0.998508\end{pmatrix}.
\]
输出层先得到
\[
 u_1=V^\top h_1+b_y\approx
 \begin{pmatrix}-0.379398\\-0.538033\end{pmatrix},\qquad
 y_1\approx\begin{pmatrix}0.406272\\0.368645\end{pmatrix}.
\]
第二步已经包含第一步的记忆：
\[
 U^\top x_2=\begin{pmatrix}1.5\\-0.4\\-2.2\end{pmatrix},\qquad
 W^\top h_1\approx\begin{pmatrix}1.668899\\-1.241174\\-0.954927\end{pmatrix},
\]
\[
 z_2\approx\begin{pmatrix}3.368899\\-1.441174\\-2.954927\end{pmatrix},\qquad
 h_2\approx\begin{pmatrix}0.997632\\-0.893934\\-0.994590\end{pmatrix}.
\]
继续使用未舍入的状态递推，四步结果为
\begin{center}
\small
\begin{tabular}{@{}cll@{}}
\toprule
$t$ & $h_t^\top$ & $y_t^\top$ \\
\midrule
1 & $(0.1974,\ 0.7163,\ {-0.9985})$ & $(0.4063,\ 0.3686)$ \\
2 & $(0.9976,\ {-0.8939},\ {-0.9946})$ & $(0.9744,\ 0.1052)$ \\
3 & $(0.9880,\ 0.2959,\ {-1.0000})$ & $(0.8620,\ 0.1765)$ \\
4 & $(0.3093,\ 0.7086,\ 0.7872)$ & $(0.5531,\ 0.6845)$ \\
\bottomrule
\end{tabular}
\end{center}
第三步最后一个隐藏分量只是四舍五入为 $-1.0000$，未舍入值约 $-0.99999725$。两个输出独立经过 sigmoid，不要求和为 $1$。

\textbf{辅助核算：}参数量为 $2\cdot3+3\cdot3+3\cdot2+3+2=26$；四个时间步共享这些参数，不乘以序列长度。以上是前向数值复算，不是训练结果，也未运行讲义中提及但未提供的 \texttt{eg8.1.ipynb}。
\end{workedsolution}

\exercisequestion{SC4001-L08-E02}{含多个时间延迟的序列回归}

\textbf{对应知识点：}\relatedtopic{SC4001-L08-T03}{长期依赖}、\relatedtopic{SC4001-L08-T04}{LSTM}。

\textbf{题目（p.~56，Example 2）。}均匀随机生成 64 条二维输入序列，每条长 16，$x(t)=(x_1(t),x_2(t))\in[0,1]^2$。构造一维目标序列：
\[
 y(t)=5x_1(t-1)x_2(t-2)-2x_1(t-7)
      +3.5x_2^2(t-5)+0.1\epsilon,
 \qquad\epsilon\sim\mathcal N(0,1).
\]
以学习率 $\alpha=0.001$ 训练 LSTM 学习输入到输出的映射，再用普通 RNN 与 GRU 比较表现。以下仅分析实验设计与依赖结构；配套 notebook 及训练结果未提供。

\begin{workedsolution}
采用 batch-first 表示时，输入形状是 $64\times16\times2$，目标形状是 $64\times16\times1$。三个轴分别为序列条数、时间长度与每步特征数；64 不是特征维数。

目标含一、二、五、七步延迟及非线性交互。以第八步为例：
\[
 y(8)=5x_1(7)x_2(6)-2x_1(1)+3.5x_2^2(3)+0.1\epsilon.
\]
因此仅看当前 $x(8)$ 不足以预测目标；模型必须从历史保留多个不同时间位置的信息，再完成乘积与平方等非线性组合。若高斯噪声独立于输入，就不能由输入精确预测，所以不应以每个样本完全拟合作为默认成功标准。

\textbf{复现边界。}\verify{当 $t\le7$ 时会引用 $t\le0$ 的历史输入；讲义未明确初始化或补历史规则，需查看配套 notebook 后确认。}图中提到的 \texttt{eg8.2a.ipynb}、\texttt{eg8.2b.ipynb} 未随本次资料提供，训练／验证划分、隐藏宽度、轮数及最终误差也不能由该页确定。本次未重跑实验，不能断言 LSTM、RNN、GRU 的实际优劣顺序。

若以后复现，需先明确上述边界，并固定数据划分与评估口径，再比较模型；这是复现建议，不是声称原讲义已经规定了这些配置。
\end{workedsolution}
